\chapter*{Введение}
\addcontentsline{toc}{chapter}{Введение}

В рамках расчетно-графической работы выполнен финальный микросервисный проект
\texttt{likes-s18}. Проект представляет собой небольшую распределенную систему, состоящую из
нескольких независимых сервисов, отдельных баз данных PostgreSQL, REST API для внешних клиентов
и gRPC-взаимодействия для внутреннего обмена данными. Такой формат работы позволяет не только
написать отдельные HTTP-обработчики, но и показать полный цикл проектирования сервиса: от
выделения доменных границ до контейнеризации и автоматической проверки в CI \cite{fastapi_docs,grpc_docs}.

Актуальность работы определяется практической необходимостью проектировать системы, которые:
\begin{itemize}
    \item можно запускать и обновлять независимо по сервисам;
    \item устойчивы к частичным отказам;
    \item удобно разворачиваются в контейнерной среде.
\end{itemize}

Цель РГР -- спроектировать и реализовать микросервисное приложение, состоящее из
\texttt{user-service}, \texttt{post-service} и \texttt{like-service}. Важной частью реализации
является сценарий, при котором \texttt{post-service} получает количество лайков по gRPC-вызову
\texttt{GetLikesCount}, а затем возвращает клиенту REST-ответ с дополнительным полем
\texttt{likes\_count}.

Для достижения цели были поставлены и решены следующие задачи:
\begin{itemize}
    \item выделение доменных границ и разделение функциональности по сервисам;
    \item разработка REST API для CRUD-операций с пользователями, постами и лайками;
    \item описание и использование gRPC-контракта на базе \texttt{.proto};
    \item организация инфраструктуры через Docker Compose и настройка healthcheck;
    \item подготовка Kubernetes-манифестов для разворачивания сервисов и баз данных;
    \item реализация CI/CD-сценария для тестирования, публикации образов и деплоя.
\end{itemize}

В отчете приведены принятые архитектурные решения, описание реализации REST и gRPC-части,
контейнеризация через Docker и Docker Compose, оркестрация в Kubernetes и настройка CI/CD-проверки
проекта \cite{week17_readme,week17_arch,kubernetes_docs,github_actions_docs}.
